\section{Reflection stability}\label{sec:upper}

Throughout this section let $X'$ be an independent copy of $X$, set
\[
Y=X-X',\qquad G\sim\Normal(0,2),
\]
and define
\[
F(z):=M_Y(z)-e^{z^2},\qquad
L:=\log(1/\Delta).
\]
For probability measures $\mu$ and $\nu$, write
\[
d_{\mathrm{TV}}(\mu,\nu):=\sup_A|\mu(A)-\nu(A)|.
\]
All constants in this section may depend on the fixed radius $\rho$, but not on $\Delta$.

\subsection{Automatic exponential integrability}

\begin{lemma}[Local reflected control supplies an exponential envelope]\label{lem:auto-envelope}
Under the hypotheses of \cref{thm:sharp-upper}, if $\Delta\le1$, then
\[
M_X(s)\le e^{s^2+\Delta}\qquad(|s|\le\rho)
\]
and
\begin{equation}\label{eq:auto-envelope}
\E e^{\rho|X|}\le M_X(\rho)+M_X(-\rho)
\le2e^{\rho^2+1}.
\end{equation}
Consequently $\E e^{\rho|Y|}$ is bounded by a constant depending only on $\rho$.
\end{lemma}

\begin{proof}
Since $\E X=0$, Jensen's inequality gives $M_X(s)\ge e^{s\E X}=1$ for every $s$ in the domain of the MGF. Hence $\log M_X(s)\ge0$. The defect bound gives
\[
\log M_X(s)+\log M_X(-s)\le s^2+\Delta,
\]
so each of the two nonnegative logarithms is at most $s^2+\Delta$. This proves the first bound. Moreover $e^{\rho|x|}\le e^{\rho x}+e^{-\rho x}$, which gives \eqref{eq:auto-envelope}. Finally,
\[
\E e^{\rho|Y|}
\le \E e^{\rho(|X|+|X'|)}
=(\E e^{\rho|X|})^2.
\]
\end{proof}

\subsection{Half-disk propagation}

\begin{lemma}[Two-constants estimate on a half-disk]\label{lem:two-constants}
Let $D_R^+=\{z\in\C:|z|<R,\ \Im z>0\}$. Suppose that $H$ is holomorphic on a neighborhood of $\overline{D_R^+}$ and satisfies
\[
|H(x)|\le a\quad (|x|\le R),
\qquad
|H(z)|\le b\quad (|z|=R,\ \Im z\ge0),
\]
where $0<a\le b$. Then for every fixed $0<\eta<1$ there is $\omega_\eta>0$, depending only on $\eta$, such that
\[
\log|H(z)|\le \omega_\eta\log a+(1-\omega_\eta)\log b,
\qquad |z|\le\eta R,\quad \Im z\ge0.
\]
The same estimate holds in the lower half-disk.
\end{lemma}

\begin{proof}
For $\delta>0$, the function $u_\delta(z):=\log(|H(z)|+\delta)$ is subharmonic. Its boundary values are at most $\log(a+\delta)$ on the diameter and at most $\log(b+\delta)$ on the semicircle. The two-constants theorem for subharmonic functions therefore bounds $u_\delta$ by the corresponding harmonic-measure combination. Under the scaling $z=Rw$, the half-disk $D_R^+$ becomes the unit half-disk, and the inner region $|z|\le\eta R$ becomes $|w|\le\eta$. The harmonic measure of the diameter is positive on this compact inner region, so its minimum is a constant $\omega_\eta>0$ depending only on $\eta$, not on $R$, $m$, or $\Delta$. Letting $\delta\downarrow0$ gives the assertion.
\end{proof}

\subsection{High-order moments of the reflected variable}

\begin{lemma}[Local defect to high-order moments]\label{lem:moment-transfer}
For every $A>0$ there is $c_0=c_0(A,\rho)>0$ such that, with
\[
m=\left\lfloor c_0\frac{L}{\log L}\right\rfloor,
\]
all sufficiently small $\Delta$ satisfy
\[
\left|\E Y^k-\E G^k\right|
\le e^{-A m\log m},
\qquad 0\le k\le2m.
\]
In particular,
\[
\E|Y|^{2m}\le (C m)^m.
\]
\end{lemma}

\begin{proof}
Set $\rho_0=\rho/2$ and put $D_{\rho_0}^{\pm}:=\{z:|z|<\rho_0,\ \pm\Im z>0\}$. By \cref{lem:auto-envelope}, $M_Y$ is holomorphic in the strip $|\Re z|<\rho$, and on the two semicircular arcs $|z|=\rho_0$,
\[
|M_Y(z)|\le \E e^{\rho_0|Y|}\le C_\rho.
\]
Thus $|F(z)|\le B_0(\rho)$ there. On the real diameter, the defect hypothesis gives
\[
|F(s)|=e^{s^2}|e^{R_X(s)}-1|\le C_0(\rho)\Delta,
\qquad |s|\le\rho_0.
\]
For sufficiently small $\Delta$, the diameter bound is no larger than the arc bound. Apply \cref{lem:two-constants} with a fixed $\eta<1$ to $D_{\rho_0}^{+}$ and $D_{\rho_0}^{-}$. For $r_0:=\eta\rho_0$, write $\omega_0:=\omega_\eta$. Thus there are constants $r_0>0$, $\omega_0>0$, and $C_1<\infty$, depending only on $\rho$, such that
\[
\sup_{|z|\le r_0}|F(z)|\le C_1\Delta^{\omega_0}.
\]
The Cauchy estimate consequently gives
\[
|F^{(k)}(0)|\le C_1 k! r_0^{-k}\Delta^{\omega_0}.
\]
For $k\le2m$, use $\log(k!)\le k\log k$ and $k\log k\le2m\log(2m)$ to obtain
\[
\begin{aligned}
\log\bigl(C_1k!r_0^{-k}\Delta^{\omega_0}\bigr)
&\le -\omega_0L+2m\log(2m)+2m|\log r_0|+O_\rho(1)\\
&\le -\omega_0L+O_\rho(c_0L),
\end{aligned}
\]
when $m=\lfloor c_0L/\log L\rfloor$ and $L$ is large. Since $m\log m\le c_0L+o(L)$, choosing $c_0$ sufficiently small, depending on $(A,\rho)$, gives $|F^{(k)}(0)|\le e^{-Am\log m}$. Since $F^{(k)}(0)=\E Y^k-\E G^k$, the moment conclusion follows. Finally,
\[
\E G^{2m}=\frac{(2m)!}{m!}\le(Cm)^m,
\]
and the negligible moment mismatch gives the stated $2m$-th moment bound.
\end{proof}

\subsection{Tail bound and truncation}

\begin{lemma}[Square-root tail lifting]\label{lem:tail-lift}
There are constants $C,c>0$ with the following property. For every sufficiently large fixed $B$, if $N=B\sqrt m$, then
\[
\Pp(|X|>N)\le C e^{-\beta_Bm},
\]
where $\beta_B\to\infty$ as $B\to\infty$. If $X_N$ denotes the conditional law $X\mid\{|X|\le N\}$, then
\[
d_{\mathrm{TV}}(\mathcal L(X),\mathcal L(X_N))\le C e^{-\beta_Bm}.
\]
Moreover
\[
|\E X_N|+|\operatorname{Var}(X_N)-1|
\le C(1+N^2)e^{-cN}.
\]
\end{lemma}

\begin{proof}
For every $p\ge1$, conditional Jensen and $\E(Y\mid X)=X$ give
\[
|X|^p=|\E(Y\mid X)|^p\le \E(|Y|^p\mid X),
\]
and hence $\E|X|^p\le\E|Y|^p$. Taking $p=2m$ and using \cref{lem:moment-transfer} gives $\|X\|_{2m}\le C\sqrt m$. Therefore, for $N=B\sqrt m$,
\[
\Pp(|X|>N)
\le \left(\frac{\|X\|_{2m}}{N}\right)^{2m}
\le C\left(\frac{C_0}{B^2}\right)^m.
\]
Thus one may take $\beta_B=\frac12\log(B^2/C_0)$ for sufficiently large $B$. If $p_N=\Pp(|X|>N)$, conditioning gives
\[
d_{\mathrm{TV}}(\mathcal L(X),\mathcal L(X_N))\le p_N.
\]
By \eqref{eq:auto-envelope}, the exponential envelope implied by the defect bound also gives, for some $c=c(\rho)>0$,
\[
\E\bigl(|X|\mathbf1_{\{|X|>N\}}\bigr)\le Ce^{-cN},
\qquad
\E\bigl(X^2\mathbf1_{\{|X|>N\}}\bigr)\le C(1+N^2)e^{-cN}.
\]
Since $\E X=0$ and $\E X^2=1$,
\[
|\E X_N|\le\frac{Ce^{-cN}}{1-p_N},
\qquad
|\E X_N^2-1|\le\frac{p_N+C(1+N^2)e^{-cN}}{1-p_N}.
\]
It follows that $|\E X_N|+|\operatorname{Var}(X_N)-1|\le C(1+N^2)e^{-cN}$ after increasing $C$.
\end{proof}

\subsection{Product approximation on a growing interval}

\begin{lemma}[Real-axis product approximation]\label{lem:real-window}
For $B$ fixed sufficiently large, there is $a_0=a_0(B,\rho)>0$ such that, with
\[
R=a_0\sqrt m
\]
and $\varphi_N$ the characteristic function of $X_N$,
\[
\sup_{|t|\le2R}
\left|\varphi_N(t)\varphi_N(-t)-e^{-t^2}\right|
\le e^{-\Gamma m}
\]
for the finite target exponent $\Gamma>4$ selected in Appendix~\ref{app:parameters}.
\end{lemma}

\begin{proof}
Use the $2m-1$ Taylor polynomial of the characteristic function of $Y=X-X'$ and the real-variable remainder
\[
\left|e^{iu}-\sum_{k=0}^{2m-1}\frac{(iu)^k}{k!}\right|
\le\frac{|u|^{2m}}{(2m)!}.
\]
Taylor's inequality and the moment comparison give
\[
\begin{aligned}
|\varphi_Y(t)-e^{-t^2}|
&\le \frac{|t|^{2m}\E|Y|^{2m}}{(2m)!}
 +\frac{|t|^{2m}\E|G|^{2m}}{(2m)!}\\
&\quad+\sum_{k=0}^{2m-1}\frac{|t|^k}{k!}
 |\E(iY)^k-\E(iG)^k|,
\end{aligned}
\]
where $G\sim\Normal(0,2)$. On $|t|\le2R$, the two Taylor-remainder terms are bounded by $(Ca_0^2)^m$ after Stirling, while the moment-mismatch term is super-exponentially smaller by \cref{lem:moment-transfer}. Choose $a_0$ small. Since $\varphi_Y(t)=\varphi_X(t)\varphi_X(-t)$ and the truncation error is bounded by \cref{lem:tail-lift}, transfer to $X_N$ gives
\[
\varphi_N(t)\varphi_N(-t)=\varphi_{X_N-X_N'}(t)=e^{-t^2}+O(e^{-\Gamma m}).
\]
The constants in the $O(\cdot)$ term are absorbed by the parameter choices in Appendix~\ref{app:parameters}.
\end{proof}

\subsection{A zero-free disk}

\begin{lemma}[Growing zero-free disk]\label{lem:zero-free}
The constants $B$ and $a_0$ can be chosen so that, for $R=a_0\sqrt m$,
\[
\varphi_N(z)\ne0,
\qquad |z|\le R.
\]
More precisely, for some $\gamma_*>a_0^2$,
\[
\sup_{|z|\le R}
\left|\varphi_N(z)\varphi_N(-z)-e^{-z^2}\right|
\le e^{-\gamma_*m}.
\]
\end{lemma}

\begin{proof}
Since $|X_N|\le N$, $\varphi_N$ is entire and
\[
|\varphi_N(z)|\le e^{N|\Im z|}.
\]
Hence, for
\[
\mathcal F_N(z):=\varphi_N(z)\varphi_N(-z)-e^{-z^2},
\]
the explicit growth bound
\[
\log|\mathcal F_N(z)|\le\log2+4NR+4R^2
\le\log2+4(Ba_0+a_0^2)m,
\qquad |z|\le2R,
\]
holds. On the real diameter, \cref{lem:real-window} gives $\log|\mathcal F_N|\le-\Gamma m$. Apply \cref{lem:two-constants} with radius $2R$ and $\eta=1/2$ in the upper and lower half-disks. With $\omega_0:=\omega_{1/2}$, this gives
\[
\log|\mathcal F_N(z)|
\le-\gamma_{\mathrm{raw}}m+O(1),
\qquad
\gamma_{\mathrm{raw}}:=\omega_0\Gamma-4(1-\omega_0)(Ba_0+a_0^2).
\]
Appendix~\ref{app:parameters} chooses $B$ and then $a_0$ so that $\gamma_{\mathrm{raw}}>a_0^2$. Choose once and for all
\[
a_0^2<\gamma_*<\gamma_{\mathrm{raw}}.
\]
For all sufficiently large $m$, the $O(1)$ term is then absorbed and
\[
|\mathcal F_N(z)|\le e^{-\gamma_*m},
\qquad |z|\le R.
\]
Since
\[
\inf_{|z|\le R}|e^{-z^2}|\ge e^{-R^2}=e^{-a_0^2m},
\]
we have
\[
\frac{|\mathcal F_N(z)|}{|e^{-z^2}|}
\le e^{-(\gamma_*-a_0^2)m}<\frac12
\]
for all sufficiently large $m$. Hence
\[
|\varphi_N(z)\varphi_N(-z)|
\ge\frac12|e^{-z^2}|>0,
\]
and therefore $\varphi_N(z)\ne0$ on the disk.
\end{proof}

\subsection{From zero-freeness to Gaussian approximation}

\begin{lemma}[From zero-freeness to Gaussian approximation]\label{lem:factor-gaussian}
Under the conclusions of \cref{lem:zero-free},
\[
\dK(X_N,\Normal(\mu_N,\sigma_N^2))\le \frac{C}{R},
\]
where $\mu_N=\E X_N$ and $\sigma_N^2=\operatorname{Var}(X_N)$. Consequently
\[
\dK(X,\Normal(0,1))\le \frac{C}{\sqrt m}.
\]
\end{lemma}

\begin{proof}
The disk $D_R:=\{z:|z|<R\}$ is simply connected, and \cref{lem:zero-free} gives $\varphi_N\ne0$ there. Hence there is a unique analytic logarithm on $D_R$ satisfying $\psi_N(0)=0$; write
\[
\psi_N(z)=\log\varphi_N(z),\qquad \psi_N(0)=0.
\]
The bounded support gives $\Re\psi_N(z)=\log|\varphi_N(z)|\le NR$. Set
\[
H(z):=NR-\psi_N(z),\qquad \Re H(z)\ge0.
\]
After scaling $z=Rw$, the Carath\'eodory coefficient theorem for functions with nonnegative real part (see, for example, Duren~\cite{Duren1983}) applied to $H(Rw)/(NR)$ gives, for $\psi_N(z)=\sum_{k\ge1}a_kz^k$,
\[
|a_k|\le2NR^{1-k}.
\]
Since $a_1=i\mu_N$ and $a_2=-\sigma_N^2/2$, for every fixed $0<\theta<1$,
\[
\left|\psi_N(t)-i\mu_Nt+\frac{\sigma_N^2t^2}{2}\right|
\le \sum_{k\ge3}2NR^{1-k}|t|^k
\le \frac{2N|t|^3}{R^2(1-|t|/R)}
\le C\frac{|t|^3}{R},
\qquad |t|\le\theta R,
\]
because $N/R=B/a_0$ is fixed. Define the cubic remainder
\[
E_N(t):=\psi_N(t)-i\mu_Nt+\frac{\sigma_N^2t^2}{2}.
\]
By \cref{lem:tail-lift}, for all sufficiently large $m$,
\[
|\mu_N|=o(R^{-1}),
\qquad
\frac12\le\sigma_N^2\le\frac32,
\qquad
|E_N(t)|\le C\frac{|t|^3}{R}.
\]
Choose $\theta$ so small that $C\theta\le1/8$. Then, on $|t|\le\theta R$,
\[
\Re\psi_N(t)=-\frac{\sigma_N^2t^2}{2}+\Re E_N(t)
\le-\frac{t^2}{4}+\frac{t^2}{8}=-\frac{t^2}{8}.
\]
Writing
\[
\varphi_N(t)=e^{i\mu_Nt-\sigma_N^2t^2/2}e^{E_N(t)},
\]
and using $|e^z-1|\le|z|e^{|z|}$ gives Gaussian damping in the difference:
\[
\left|\varphi_N(t)-e^{i\mu_Nt-\sigma_N^2t^2/2}\right|
\le C\frac{|t|^3}{R}e^{-ct^2}.
\]
Let $g_N$ be the density of $\Normal(\mu_N,\sigma_N^2)$. The Esseen smoothing inequality in the present normalization (Esseen~\cite{Esseen1945}) is
\[
\begin{aligned}
\dK(X_N,\Normal(\mu_N,\sigma_N^2))
&\le \frac1\pi\int_{-T}^{T}
\left|\frac{\varphi_N(t)-e^{i\mu_Nt-\sigma_N^2t^2/2}}{t}\right|\,dt
 +\frac{24\|g_N\|_\infty}{\pi T}\\
&= \frac1\pi\int_{-T}^{T}
\left|\frac{\varphi_N(t)-e^{i\mu_Nt-\sigma_N^2t^2/2}}{t}\right|\,dt
 +O(T^{-1}),
\end{aligned}
\]
where $\|g_N\|_\infty=(\sqrt{2\pi}\sigma_N)^{-1}$ and the last form uses $\sigma_N^2\ge1/2$. With $T=\theta R$, the integral is bounded by
\[
\frac{C}{R}\int_{-T}^{T}t^2e^{-ct^2}\,dt
\le\frac{C}{R}\int_\R t^2e^{-ct^2}\,dt
 =O(R^{-1}),
\]
and the cutoff term is also $O(R^{-1})$.

It remains to compare the nearby Gaussian with the standard one. For $\sigma$ in a fixed compact subset of $(0,\infty)$,
\begin{equation}\label{eq:gaussian-parameter-kolmogorov}
\dK(\Normal(\mu,\sigma^2),\Normal(0,1))
\le C\bigl(|\mu|+|\sigma-1|\bigr).
\end{equation}
Indeed, the translation term is bounded by $C|\mu|$ using the bounded Gaussian density, while the mean-value theorem in $\sigma$ gives
\[
\sup_x|\Phi(x/\sigma)-\Phi(x)|
\le C|\sigma-1|,
\]
because $\sup_y|y|\phi(y)<\infty$. Since $\sigma_N^2\in[1/2,3/2]$, also $|\sigma_N-1|\le C|\sigma_N^2-1|$. Thus \cref{lem:tail-lift} and \eqref{eq:gaussian-parameter-kolmogorov} show that the Gaussian parameter shift is $o(R^{-1})$. Finally,
\[
\dK(X,X_N)\le d_{\mathrm{TV}}(\mathcal L(X),\mathcal L(X_N))=o(R^{-1}),
\]
so the triangle inequality yields the consequence for $X$.
\end{proof}

\begin{proof}[Proof of \cref{thm:sharp-upper}]
First apply \cref{lem:auto-envelope}; this is exactly where the exponential envelope formerly imposed as a separate hypothesis is obtained from the reflected defect. Then apply \cref{lem:moment-transfer,lem:tail-lift,lem:real-window,lem:zero-free,lem:factor-gaussian}. Since
\[
m\asymp_\rho\frac{L}{\log L},\qquad L=\log(1/\Delta),
\]
we have
\[
\frac1{\sqrt m}
\asymp_\rho
\sqrt{\frac{\log L}{L}},
\]
which is the claimed rate.
\end{proof}
